%% Question bank chapter file - generated by bank_tools.py build (do not edit while the bank is still being built:
%% it is regenerated from the per-test data in Question Bank\data). Every question is one bankq block:
%%   \begin{bankq}[id=..., source={...}, paper=.., qno=.., type=..., subtopic={...}, answer={...}, answer_src=key|solved,
%%                 working={...}, difficulty=1-10, diff_parts={K S I C F}, diff_note={...},
%%                 group=<passage id or empty>, also_in={duplicate copies in other tests}]
%% difficulty = min(10, K+S+I+C+F-1): K concepts 1-3, S steps 1-4, I insight 0-2, C computation 0-2, F format 0-1
%%   ...question LaTeX, options, figures (\bfig / \bfigi -> img/)...
%%   \end{bankq}
%% Passages shared by several questions are bankpassage blocks; their questions carry group=<passage id>.
\documentclass[10pt]{article}
\usepackage[a4paper,margin=1.8cm]{geometry}
\usepackage[T1]{fontenc}
\usepackage{amsmath,amssymb,mathtools}
\usepackage{graphicx}
\usepackage[export]{adjustbox}
\usepackage{xcolor,needspace,keyval}
\usepackage{enumitem}
\usepackage{booktabs,array,tabularx}
\usepackage[hidelinks]{hyperref}
\setlength{\parindent}{0pt}
\setlength{\parskip}{3pt}
\makeatletter
\define@key{bq}{id}{\def\bq@id{#1}}\define@key{bq}{source}{\def\bq@src{#1}}\define@key{bq}{paper}{\def\bq@paper{#1}}
\define@key{bq}{qno}{\def\bq@qno{#1}}\define@key{bq}{type}{\def\bq@type{#1}}\define@key{bq}{subtopic}{\def\bq@sub{#1}}
\define@key{bq}{answer}{\def\bq@ans{#1}}\define@key{bq}{answer_src}{\def\bq@asrc{#1}}\define@key{bq}{working}{\def\bq@work{#1}}
\define@key{bq}{group}{\def\bq@group{#1}}\define@key{bq}{also_in}{\def\bq@also{#1}}
\define@key{bq}{difficulty}{\def\bq@diff{#1}}\define@key{bq}{diff_parts}{\def\bq@dparts{#1}}\define@key{bq}{diff_note}{\def\bq@dnote{#1}}
\newenvironment{bankq}[1][]{\def\bq@work{}\def\bq@also{}\def\bq@group{}\def\bq@diff{?}\def\bq@dparts{}\def\bq@dnote{}\setkeys{bq}{#1}%
  \par\needspace{4\baselineskip}\medskip\noindent{\small\textbf{\bq@id}\hfill\color{gray}[\bq@src, Paper \bq@paper, Q\bq@qno]
  \bq@type\ $\cdot$ \bq@sub\ $\cdot$ \textbf{D\bq@diff}}\par\nopagebreak\relax}
  {\par{\small\color{teal}Answer: \textbf{\bq@ans}\ifx\bq@work\@empty\else\ (solved: \bq@work)\fi
  \ifx\bq@also\@empty\else\par\color{gray}Also in: \bq@also\fi}\par}
\define@key{bp}{id}{\def\bp@id{#1}}\define@key{bp}{source}{\def\bp@src{#1}}\define@key{bp}{questions}{\def\bp@qs{#1}}
\newenvironment{bankpassage}[1][]{\setkeys{bp}{#1}\par\needspace{4\baselineskip}\bigskip\noindent
  {\small\textbf{Passage \bp@id}\hfill\color{gray}[\bp@src; questions \bp@qs]}\par\nopagebreak\relax\itshape}{\par}
\makeatother
\newcommand{\bfig}[2]{\par\smallskip{\centering\includegraphics[width=#1pt,max width=0.62\linewidth,max height=0.32\textheight]{img/#2}\par}\smallskip}
\newcommand{\bfigi}[2]{\includegraphics[width=#1pt,max width=\linewidth,max height=0.32\textheight,valign=c]{img/#2}}
\newcommand{\bmissing}[1]{\textit{[figure missing in the source conversion: \detokenize{#1}]}}
\newenvironment{qoptions}{\begin{itemize}[leftmargin=1.8em,itemsep=1pt,topsep=2pt,label={}]}{\end{itemize}}
\newcommand{\opt}[2]{\item[(#1)] #2}
\providecommand{\fwstart}{}\providecommand{\fwend}{}\providecommand{\qfullsplit}{}
\newcolumntype{L}[1]{>{\raggedright\arraybackslash}p{#1}}
\newcolumntype{C}[1]{>{\centering\arraybackslash}p{#1}}
\newcolumntype{R}[1]{>{\raggedleft\arraybackslash}p{#1}}
\providecommand{\sectionhdr}[1]{}\providecommand{\subsectionhdr}[1]{}


\begin{document}
{\Large\bfseries Parabola --- Question Bank}\par\smallskip
{\small Mathematics chapter M17. Questions from all converted test papers, grouped by source series. Answers: printed key where it matches the paper, otherwise solved (answer\_src=solved, with working). Exact repeats across papers appear once, with also\_in listing the other copies.}\par\medskip\hrule

\section*{Score-I papers}
\begin{bankq}[id=SI12-P1-MATH-S1-Q02, source={Score-I - Test 12 (09 Mar 2022)}, paper=1, qno=2, type=multi, subtopic={Tangents}, answer={A,B,C,D}, answer_src=key, difficulty=7, diff_parts={K2 S3 I1 C1 F1}, diff_note={concepts: tangent properties of a parabola, focal properties | steps: intersection of the two tangents; tangent at the vertex and axis direction; vertex; focus; latus rectum; verify each option | I: foot of the perpendicular from the focus lies on the vertex tangent}, also_in={Nurture (General) - Test 04 (11 Feb 2024) P1 Q8}]
If equation of tangent at $P$, $Q$ and vertex $A$ of a parabola are $3x+4y-7=0$, $2x+3y-10=0$ and $x-y=0$ respectively, then
\begin{qoptions}
\opt{A}{focus is $(4,5)$}
\opt{B}{length of latus rectum is $2\sqrt2$}
\opt{C}{axis is $x+y-9=0$}
\opt{D}{vertex is $\left(\dfrac92,\dfrac92\right)$}
\end{qoptions}
\end{bankq}

\section*{Nurture SRG papers}
\begin{bankq}[id=NSRG03-P1-MATH-S2-Q11, source={Nurture SRG - Test 03 (04 Feb 2024)}, paper=1, qno=11, type=multi, subtopic={Standard forms}, answer={A,C}, answer_src=solved, working={Parabola with focus $(1,0)$ and directrix the tangent $x-2y+4=0$; latus rectum $2\sqrt5$}, difficulty=8, diff_parts={K1 S3 I2 C2 F1}, diff_note={concepts: locus of the centre of a circle touching a tangent and passing through the focus | steps: tangent at the point; definition as a parabola; directrix; focus; axis; latus rectum}]
Let a variable circle touches a line, which is tangent to the parabola $y^2=4x$ at $(4,4)$ and variable circle passes through focus of the parabola $y^2=4x$. If locus of its centre is conic `C' then
\begin{qoptions}
\opt{A}{Directrix of the conic `C' is $x-2y+4=0$}
\opt{B}{Focus of the conic `C' is $(4,0)$}
\opt{C}{Latus rectum of conic `C' is $2\sqrt5$}
\opt{D}{Axis of the conic `C' is $x+2y-1=0$}
\end{qoptions}
\end{bankq}
\begin{bankq}[id=NSRG03-P1-MATH-S2-Q12, source={Nurture SRG - Test 03 (04 Feb 2024)}, paper=1, qno=12, type=multi, subtopic={Standard forms}, answer={B,C}, answer_src=solved, working={$P(4,4)$; circle with centre $(2.5,2)$ and radius $2.5$ touches $x=0$; its chord on the normal is $\sqrt5$}, difficulty=8, diff_parts={K1 S3 I2 C2 F1}, diff_note={concepts: circle on a focal radius of a parabola | steps: point of the parabola; focus; circle; tangency to the vertex tangent; normal intercept; radius}]
Let $P(t^2,2t)$ be the point on the parabola $y^2=4x$ other than its vertex $A$ whose abscissa and ordinate are same. A circle $C$ is drawn with focal radius of the point $P$ as diameter, then which of the following hold(s) good ?
\begin{qoptions}
\opt{A}{Circle $C$ touches the directrix of the parabola.}
\opt{B}{Circle $C$ touches the tangent at $A$ of the parabola.}
\opt{C}{Length of the intercept made by the circle $C$ on the normal drawn at $P$ of the parabola is $\sqrt5$}
\opt{D}{Radius of the circle $C$ is 5.}
\end{qoptions}
\end{bankq}
\begin{bankq}[id=NSRG03-P1-MATH-S3-Q02, source={Nurture SRG - Test 03 (04 Feb 2024)}, paper=1, qno=2, type=numerical, subtopic={Chords}, answer={200}, answer_src=solved, working={Projection condition gives $(4h-k^2)(k+2)^2=200$}, difficulty=7, diff_parts={K1 S2 I2 C2 F1}, diff_note={concepts: locus of midpoint of a chord with fixed projection | steps: midpoint in terms of the ordinates; projection condition; relation of the sums; eliminate}]
$PQ$ is a variable chord of the parabola $y^2=4x$ such that its projection with a line with gradient unity has a magnitude equal to 10. If the locus of the middle point of $PQ$ has the equation $\left(y^2-4x\right)(y+2)^2+\lambda=0$ where $\lambda\in N$, then find $\lambda$.
\end{bankq}
\begin{bankq}[id=NSRG03-P2-MATH-S1-Q03, source={Nurture SRG - Test 03 (04 Feb 2024)}, paper=2, qno=3, type=single, subtopic={Chords}, answer={D}, answer_src=solved, working={$(4h-k^2)(k^2+4)=4\ell^2$}, difficulty=6, diff_parts={K1 S2 I2 C2 F0}, diff_note={concepts: locus of midpoints of constant length chords of a parabola | steps: midpoint relations; length condition; eliminate; locus}]
The locus of the middle points of chords of the parabola $y^2=4x$, which are of constant length `$2\ell$' is :
\begin{qoptions}
\opt{A}{$\left(4x+y^2\right)\left(y^2-4\right)=4\ell^2$}
\opt{B}{$\left(4y+x^2\right)\left(x^2-4\right)=4\ell^2$}
\opt{C}{$\left(4y-x^2\right)\left(x^2+4\right)=4\ell^2$}
\opt{D}{$\left(4x-y^2\right)\left(y^2+4\right)=4\ell^2$}
\end{qoptions}
\end{bankq}
\begin{bankq}[id=NSRG03-P2-MATH-S2-Q08, source={Nurture SRG - Test 03 (04 Feb 2024)}, paper=2, qno=8, type=multi, subtopic={Normals}, answer={B,C}, answer_src=solved, working={Three normals need the cubic $am^3+(2a-\sin\theta)m+\cos\theta=0$ to have three real roots: $a\in(-\tfrac12,0)\cup(0,\tfrac12)$}, difficulty=7, diff_parts={K1 S2 I2 C2 F1}, diff_note={concepts: three normals from a point on a unit circle to a parabola | steps: equation of the normals; cubic in the slope; condition for three real roots; range of a}]
From a point $(\sin\theta,\cos\theta)$ if three normals can be drawn to the parabola $y^2=4ax$ then the value of `$a$' belong to
\begin{qoptions}
\opt{A}{$\left(\dfrac12,1\right)$}
\opt{B}{$\left(-\dfrac12,0\right)$}
\opt{C}{$\left(0,\dfrac12\right)$}
\opt{D}{$(1,\infty)$}
\end{qoptions}
\end{bankq}
\begin{bankq}[id=NSRG03-P2-MATH-S2-Q12, source={Nurture SRG - Test 03 (04 Feb 2024)}, paper=2, qno=12, type=multi, subtopic={Chords}, answer={C}, answer_src=solved, working={Midpoint ordinate satisfies $mk^2-(4m+2)k+8m=0$; two genuine distinct chords only for $m=\tfrac12$ (for $m=0$ and $m=1$ only one genuine chord; the printed key lists B,C,D)}, difficulty=7, diff_parts={K1 S2 I2 C2 F1}, diff_note={concepts: chords from a point of a parabola bisected by a line | steps: midpoint relation; quadratic in the ordinate; two distinct genuine chords; range of the slope}]
$y=mx$ bisects two distinct chords drawn from $(4,4)$ on $y^2=4x$ if
\begin{qoptions}
\opt{A}{$m=-1/2$}
\opt{B}{$m=0$}
\opt{C}{$m=1/2$}
\opt{D}{$m=1$}
\end{qoptions}
\end{bankq}
\begin{bankq}[id=NSRG03-P2-MATH-S3-Q03, source={Nurture SRG - Test 03 (04 Feb 2024)}, paper=2, qno=3, type=numerical, subtopic={Tangents}, answer={5}, answer_src=solved, working={Centre $(-9,14)$, radius $5\sqrt{10}$}, difficulty=7, diff_parts={K1 S2 I2 C2 F1}, diff_note={concepts: circle through the focus touching a parabola at a point | steps: tangent and normal; centre on the normal; equal distance to the focus; radius}]
The radius of circle which passes through the focus of parabola $x^2=4y$ and touches it at point $(6,9)$ is $k\sqrt{10}$, then $k=$
\end{bankq}

\section*{TOAS papers}
\begin{bankq}[id=TOAS02-P1-MATH-S2-Q07, source={TOAS - Test 02 (13 Feb 2022)}, paper=1, qno=7, type=multi, subtopic={Tangents}, answer={B,C,D}, answer_src=key, difficulty=5, diff_parts={K2 S2 I1 C0 F1}, diff_note={concepts: tangent to a parabola, chord of contact of a circle | steps: tangent at (x1, y1); chord of contact from a point on it; fixed point (x2, y2); verify the relations | I: chords of contact form a pencil of lines}]
Tangent is drawn at any point $(x_1,y_1)$ other than the vertex, on parabola $y^2=4ax$. If tangents are drawn from any point on this tangent to the circle $x^2+y^2=a^2$ such that all the chords of contact pass through a fixed point $(x_2,y_2)$ then
\begin{qoptions}
\opt{A}{$x_1,a,x_2$ are in GP}
\opt{B}{$\dfrac{y_1}{2},a,y_2$ are in GP}
\opt{C}{$-4,\dfrac{y_1}{y_2},\dfrac{x_1}{x_2}$ are in GP}
\opt{D}{$x_1x_2+y_1y_2=a^2$}
\end{qoptions}
\end{bankq}
\begin{bankq}[id=TOAS02-P1-MATH-S2-Q08, source={TOAS - Test 02 (13 Feb 2022)}, paper=1, qno=8, type=multi, subtopic={Chords}, answer={B,D}, answer_src=key, difficulty=3, diff_parts={K1 S2 I0 C0 F1}, diff_note={concepts: focal chord of a parabola | steps: substitute the line in the parabola; chord length for m = 1; normals at the ends; angle between them}]
A straight line $y=m(x-a)$, $m\ne0$ meets a parabola $y^2=4ax$ at two points $A$ and $B$. Then which of the following option(s) is/are correct
\begin{qoptions}
\opt{A}{Length $AB=4a$ when $m=1$}
\opt{B}{Length $AB=8a$ when $m=1$}
\opt{C}{Angle between normals drawn at $A$ and $B$ is $45^\circ$ when $m=1$}
\opt{D}{Angle between normals drawn at $A$ and $B$ is $90^\circ$ when $m=2$}
\end{qoptions}
\end{bankq}
\begin{bankq}[id=TOAS12-P1-MATH-S1-Q07, source={TOAS - Test 12 (10 Jul 2022)}, paper=1, qno=7, type=multi, subtopic={Normals}, answer={A,D}, answer_src=key, difficulty=4, diff_parts={K1 S2 I1 C0 F1}, diff_note={concepts: normals and tangents of a parabola | steps: condition for a normal through the chord; intersection of tangents; locus and integer points | I: relation between the end parameters of a normal chord}]
Let line $L$ is normal to parabola $y^2=4x$ and line $L$ intersect the parabola at $A$ and $B$. If point $P(\alpha,\beta)$ is point of intersection of tangents at $A$ and $B$ then which of the following is/are correct?
\begin{qoptions}
\opt{A}{Locus of point $P$ is symmetric w.r.t. $x$ axis}
\opt{B}{Locus of point $P$ is symmetric w.r.t. $y$ axis}
\opt{C}{If $\alpha,\beta\in I$ then total number of possible ordered pair $(\alpha,\beta)$ is 2}
\opt{D}{If $\alpha,\beta\in I$ then total number of possible ordered pair $(\alpha,\beta)$ is 4}
\end{qoptions}
\end{bankq}
\begin{bankq}[id=TOAS13-P1-MATH-S3-Q01, source={TOAS - Test 13 (24 Jul 2022)}, paper=1, qno=1, type=numerical, subtopic={Standard forms}, answer={0.25}, answer_src=key, difficulty=6, diff_parts={K2 S2 I1 C1 F1}, diff_note={concepts: pencil of conics through four points, parabola condition | steps: line pairs through the points give the family; h\^{}2 = ab fixes the parameter; the quadratic part is a perfect square; axis direction slope}]
Let parabola $P$ passes through the points $(0,1)$, $(2,0)$, $(0,2)$ and $(2,2)$. If axis of symmetry of the parabola has slope $m$, then $|m|$ is equal to \underline{\hspace{1.5cm}}.
\end{bankq}
\begin{bankq}[id=TOAS13-P2-MATH-S1-Q19, source={TOAS - Test 13 (24 Jul 2022)}, paper=2, qno=19, type=single, subtopic={Normals}, answer={A}, answer_src=key, difficulty=4, diff_parts={K2 S2 I1 C0 F0}, diff_note={concepts: normals to a parabola and circle on a diameter, power of a point | steps: normal condition gives m\^{}3 - 7m + 6 = 0; parameters of P and Q give T; circle on TR as diameter; power of (-1,-1)}]
Tangents and normals at points $P$ and $Q$ to the parabola $y^2=4x$ intersect at point $T$ and point $R$ (9, 6) respectively. Then find the length of tangent drawn from $(-1,-1)$ to the circle circumscribing the quadrilateral PTQR.
\begin{qoptions}
\opt{A}{4}
\opt{B}{8}
\opt{C}{3}
\opt{D}{1}
\end{qoptions}
\end{bankq}
\begin{bankq}[id=TOAS19-P1-MATH-S3-Q03, source={TOAS - Test 19 (05 Feb 2023)}, paper=1, qno=3, type=numerical, subtopic={Normals}, answer={1.50}, answer_src=key, difficulty=6, diff_parts={K1 S2 I2 C1 F1}, diff_note={concepts: common normal of two parabolas | steps: normal of the first parabola; condition for a common normal; range of k}, also_in={TNAS - Test 01 (05 Feb 2023) P2 Q3}]
If $y^2=4x$ and $y^2=(x-k)$ have a common normal other than $x$-axis, then the range of values of $k$ is $(p,\infty)$ ; where $p$ is equal to
\end{bankq}

\section*{TOPS papers}
\begin{bankq}[id=TOPS04-P1-MATH-S1-Q03, source={TOPS - Test 04 (18 Dec 2022)}, paper=1, qno=3, type=multi, subtopic={Tangents}, answer={A,C}, answer_src=key, difficulty=9, diff_parts={K2 S3 I2 C2 F1}, diff_note={concepts: parabola, tangents with points of contact | steps: use tangent pairs and points of contact; focus and directrix from the tangent properties; axis; vertex; check}]
$x+y=2$ and $x-y=2$ are tangents on a parabola at $(1,1)$ and $(4,2)$ respectively. Which of the following is/are correct
\begin{qoptions}
\opt{A}{Equation of directrix is $x+3y=2$}
\opt{B}{Equation of axis is $3x-y=5$}
\opt{C}{Focus of the parabola is at $\left(\dfrac85,\dfrac65\right)$}
\opt{D}{Vertex of the parabola is at $\left(\dfrac{33}{20},\dfrac{13}{20}\right)$}
\end{qoptions}
\end{bankq}
\begin{bankq}[id=TOPS04-P1-MATH-S1-Q04, source={TOPS - Test 04 (18 Dec 2022)}, paper=1, qno=4, type=multi, subtopic={Chords}, answer={A,B}, answer_src=key, difficulty=8, diff_parts={K2 S2 I2 C2 F1}, diff_note={concepts: parabolas on a common axis, locus | steps: equations of the two parabolas; intercept on a line; midpoint locus; cases of the latus rectums}]
Let two parabolas have a common axis where focus of each being exterior to the other and let $\ell_1,\ell_2$ be their latus rectums then the locus of the mid points of the intercepts between the parabolas made on the lines parallel to the common axis is a
\begin{qoptions}
\opt{A}{Straight line if $\ell_1=\ell_2\ \forall\ \ell_1,\ell_2\in R^+$}
\opt{B}{Parabola if $\ell_1\neq\ell_2\ \forall\ \ell_1,\ell_2\in R^+$}
\opt{C}{Parabola $\forall\ \ell_1,\ell_2\in R^+$}
\opt{D}{Hyperbola if $\ell_1=2\ell_2\ \forall\ \ell_1,\ell_2\in R^+$}
\end{qoptions}
\end{bankq}
\begin{bankq}[id=TOPS04-P2-MATH-S1-Q08, source={TOPS - Test 04 (18 Dec 2022)}, paper=2, qno=8, type=single, subtopic={Chords}, answer={B}, answer_src=key, difficulty=4, diff_parts={K1 S2 I1 C1 F0}, diff_note={concepts: midpoint locus of chords | steps: chord with given slope; midpoint relation; locus}]
The locus of middle point of chord of parabola $y^2=4x$ which is parallel to the line $y=5x$ is
\begin{qoptions}
\opt{A}{$x^2+y^2=7$}
\opt{B}{$y=\dfrac25$}
\opt{C}{$x=\dfrac32$}
\opt{D}{$x+y=\dfrac56$}
\end{qoptions}
\end{bankq}
\begin{bankq}[id=TOPS04-P2-MATH-S1-Q15, source={TOPS - Test 04 (18 Dec 2022)}, paper=2, qno=15, type=single, subtopic={Tangents}, answer={B}, answer_src=key, difficulty=4, diff_parts={K1 S2 I1 C1 F0}, diff_note={concepts: normal of parabola | steps: normal condition; second point; expression}]
Let the normal of the parabola $y^2=4ax$ at the point $P(at_1^2,2at_1)$, intersects the curve again at $Q(at_2^2,2at_2)$, then $t_2+t_1+\dfrac2{t_1}$ is
\begin{qoptions}
\opt{A}{$-2$}
\opt{B}{0}
\opt{C}{1}
\opt{D}{$-3$}
\end{qoptions}
\end{bankq}
\begin{bankq}[id=TOPS04-P2-MATH-S2-Q03, source={TOPS - Test 04 (18 Dec 2022)}, paper=2, qno=3, type=numerical, subtopic={Standard forms}, answer={17}, answer_src=solved, working={$\lambda_1=8$, $\lambda_2=4$, $\lambda_3=\sqrt{16+9}=5$, sum $17$}, difficulty=5, diff_parts={K1 S2 I1 C1 F1}, diff_note={concepts: parabola quantities | steps: latus rectum; focus directrix distance; focal distance}]
For the parabola $y^2=8x$, if length of latus ractum is $\lambda_1$, distance between focus and directrix is $\lambda_2$ and distance between focus and $(6,3)$ is $\lambda_3$, then $(\lambda_1+\lambda_2+\lambda_3)$ is
\end{bankq}
\begin{bankq}[id=TOPS05-P1-MATH-S1-Q02, source={TOPS - Test 05 (14 Jan 2023)}, paper=1, qno=2, type=multi, subtopic={Normals}, answer={A,C}, answer_src=key, difficulty=9, diff_parts={K3 S2 I2 C2 F1}, diff_note={concepts: parabola, normals, parameters | steps: normal condition in parameters; sum of the parameters; centroid; tangent triangle area formula}]
Normals are drawn from the point $P(15,12)$ to the parabola $y^2=4x$. If the feet of the normals form a $\Delta ABC$, then which of the following is/are true
\begin{qoptions}
\opt{A}{Centroid of $\Delta ABC$ is $\left(\dfrac{26}3,0\right)$}
\opt{B}{Centroid of $\Delta ABC$ is $\left(\dfrac{22}3,0\right)$}
\opt{C}{Area of triangle formed by tangents to the parabola at $A$, $B$ and $C$ is 35 sq. units}
\opt{D}{Area of triangle formed by tangents to the parabola at $A$, $B$ and $C$ is 140 sq. units}
\end{qoptions}
\end{bankq}
\begin{bankq}[id=TOPS05-P1-MATH-S1-Q06, source={TOPS - Test 05 (14 Jan 2023)}, paper=1, qno=6, type=multi, subtopic={Chords}, answer={A,D}, answer_src=key, difficulty=8, diff_parts={K3 S2 I2 C1 F1}, diff_note={concepts: focal chord, parabola, means | steps: parametric line through the point; the product of the lengths; mean relations; check}]
If a variable straight line passes through the point $P(2,0)$, intersects the curve $y^2=8x$ at the points $Q$, $R$ (let $PQ=\ell_1$ \& $PR=\ell_2$). Also given that $G$, $A$ represents geometric and arithmetic means of $\ell_1$ and $\ell_2$ respectively, then correct option(s) is/are
\begin{qoptions}
\opt{A}{$\ell_1,\ell_2\ge2$}
\opt{B}{$\ell_1,\ell_2<2$}
\opt{C}{$G\propto A$}
\opt{D}{$G\propto\sqrt A$}
\end{qoptions}
\end{bankq}
\begin{bankq}[id=TOPS05-P1-MATH-S3-Q07, source={TOPS - Test 05 (14 Jan 2023)}, paper=1, qno=7, type=numerical, subtopic={Normals}, answer={2.665}, answer_src=key, difficulty=10, diff_parts={K3 S3 I2 C2 F1}, diff_note={concepts: normals of a parabola, locus, right angle | steps: parameters of the chord; condition from the right angle; point of intersection of normals; eliminate; locus; ratio}]
The locus of point of intersection of the normals to the parabola $y^2=4ax$ drawn at the extremities of a chord subtending an angle of $90^\circ$ at origin is $y^2=ab(x-ac)$ then $\left(\dfrac bc\right)=$
\end{bankq}
\begin{bankq}[id=TOPS05-P2-MATH-S2-Q07, source={TOPS - Test 05 (14 Jan 2023)}, paper=2, qno=7, type=multi, subtopic={Tangents}, answer={C,D}, answer_src=key, difficulty=9, diff_parts={K2 S3 I2 C2 F1}, diff_note={concepts: parabola tangents, focus and directrix | steps: reflection property; foot of the directrix; focus; directrix; latus rectum; check}]
A parabola touches straight lines $x-2y+13=0$ and $2x+y+1=0$ at points $P(17,15)$ and $Q(2,-5)$ respectively then :
\begin{qoptions}
\opt{A}{Focus of parabola is $\left(\dfrac72,-3\right)$}
\opt{B}{Directrix of parabola is $x+y-2=0$}
\opt{C}{Focus of parabola is $(5,-1)$}
\opt{D}{Length of latus rectum of parabola is 16}
\end{qoptions}
\end{bankq}

\section*{TNAS papers}
\begin{bankq}[id=TNAS07-P1-MATH-S1-Q14, source={TNAS - Test 07 (20 Aug 2023)}, paper=1, qno=14, type=single, subtopic={Tangents}, answer={D}, answer_src=key, difficulty=5, diff_parts={K1 S2 I2 C1 F0}, diff_note={concepts: locus of the foot of the perpendicular on tangents | steps: equation of tangent; perpendicular from the origin; eliminate the slope; simplify}]
The locus of foot of perpendicular drawn from origin upon any tangent of parabola $y^2=8(x-1)$ is
\begin{qoptions}
\opt{A}{$x^3+xy^2+2y^2+x^2=0$}
\opt{B}{$x^3-xy^2+2y^2-x^2=0$}
\opt{C}{$x^3-xy^2+2y^2+x^2=0$}
\opt{D}{$x^3+xy^2+2y^2-x^2=0$}
\end{qoptions}
\end{bankq}
\begin{bankq}[id=TNAS07-P1-MATH-S1-Q16, source={TNAS - Test 07 (20 Aug 2023)}, paper=1, qno=16, type=single, subtopic={Tangents}, answer={C}, answer_src=key, difficulty=6, diff_parts={K1 S2 I2 C2 F0}, diff_note={concepts: two parabolas touching at a point | steps: point condition; equal slopes; solve for the constants; evaluate}]
Two parabolas $x=cy^2+ay+b$ and $x=cy-y^2$ touch each other at the point $(0,1)$, then the value of $2a+b-c$ is-
\begin{qoptions}
\opt{A}{$-1$}
\opt{B}{$-7$}
\opt{C}{$-5$}
\opt{D}{1}
\end{qoptions}
\end{bankq}
\begin{bankq}[id=TNAS07-P2-MATH-S1-Q01, source={TNAS - Test 07 (20 Aug 2023)}, paper=2, qno=1, type=multi, subtopic={Standard forms}, answer={A,B,C,D}, answer_src=key, difficulty=8, diff_parts={K1 S3 I2 C2 F1}, diff_note={concepts: parabola from the focus and the directrix with a condition | steps: directrix coefficients; use the curve condition; focus on the axis; equal distances; tangent condition; check options}]
If point $(4,4)$ lies on parabola whose focus lies on $x$-axis and directrix is $\ell x+my=1$ such that $(\ell,m)$ lies on curve $x^2+y^2-32xy+8x+8y-1=0$ then :
\begin{qoptions}
\opt{A}{co-ordinates of focus may be $(5,0)$}
\opt{B}{co-ordinates of focus may be $(3,0)$}
\opt{C}{If tangent at $(4,4)$ passes through origin, value of $m$ may be $\dfrac13$}
\opt{D}{If tangent at $(4,4)$ passes through origin, equation of axis may be $x-4y-3=0$}
\end{qoptions}
\end{bankq}

\section*{Nurture (General) papers}
\begin{bankq}[id=NUR01-P1-MATH-S2-Q01, source={Nurture (General) - Test 01 (12 Jun 2021)}, paper=1, qno=1, type=single, subtopic={Tangents}, answer={C}, answer_src=solved, working={a=1/2 gives the parabola sqrt x+sqrt y=1 with tangent x+y=1/2; the hyperbola with x+y=4; line x+y=1; circle with x+y=sqrt2}, difficulty=6, diff_parts={K2 S2 I2 C1 F0}, diff_note={concepts: curve x\^{}a+y\^{}a=1 table match, tangents | steps: identify each curve; tangent of slope minus one; find the correct match}]
Which of the following options is the only \textbf{CORRECT} combination?
\begin{qoptions}
\opt{A}{(III) (iii) (P)}
\opt{B}{(IV) (i) (S)}
\opt{C}{(II) (iv) (P)}
\opt{D}{(I) (ii) (R)}
\end{qoptions}
\end{bankq}
\begin{bankq}[id=NUR01-P1-MATH-S2-Q02, source={Nurture (General) - Test 01 (12 Jun 2021)}, paper=1, qno=2, type=single, subtopic={Tangents}, answer={C}, answer_src=solved, working={(III),(iii),(P) is wrong; the line x+y=1 is tangent to itself, option Q}, difficulty=6, diff_parts={K2 S2 I2 C1 F0}, diff_note={concepts: curve x\^{}a+y\^{}a=1 table incorrect match, tangents | steps: identify each curve; tangent of slope minus one; find the wrong match}]
Which of the following options is the only \textbf{INCORRECT} combination ?
\begin{qoptions}
\opt{A}{(III) (iii) (Q)}
\opt{B}{(I) (ii) (S)}
\opt{C}{(III) (iii) (P)}
\opt{D}{(II) (iv) (P)}
\end{qoptions}
\end{bankq}
\begin{bankq}[id=NUR01-P1-MATH-S2-Q03, source={Nurture (General) - Test 01 (12 Jun 2021)}, paper=1, qno=3, type=single, subtopic={Tangents}, answer={D}, answer_src=solved, working={The circle x\^{}2+y\^{}2=1 has the tangent x+y=sqrt2}, difficulty=6, diff_parts={K2 S2 I2 C1 F0}, diff_note={concepts: curve x\^{}a+y\^{}a=1 table match, tangents | steps: identify each curve; tangent of slope minus one; find the correct match}]
Which of the following options is the only \textbf{CORRECT} combination?
\begin{qoptions}
\opt{A}{(II) (iii) (P)}
\opt{B}{(I) (ii) (R)}
\opt{C}{(IV) (i) (Q)}
\opt{D}{(IV) (i) (R)}
\end{qoptions}
\end{bankq}
\begin{bankq}[id=NUR01-P2-MATH-S1-Q06, source={Nurture (General) - Test 01 (12 Jun 2021)}, paper=2, qno=6, type=single, subtopic={Standard forms}, answer={A}, answer_src=solved, working={Directrix through the reflections of S; distance 3/sqrt5; latus rectum 6/sqrt5}, difficulty=7, diff_parts={K2 S2 I2 C2 F0}, diff_note={concepts: parabola from two tangents and a focus, standard forms | steps: perpendicular tangents; directrix; distance; latus rectum}]
Let $x+y=0$ and $x-y=0$ are tangents drawn to a parabola whose focus is at $S(1,2)$ then length of latus rectum of parabola is
\begin{qoptions}
\opt{A}{$\dfrac6{\sqrt5}$}
\opt{B}{$\dfrac4{\sqrt5}$}
\opt{C}{$\dfrac2{\sqrt5}$}
\opt{D}{$\dfrac1{\sqrt5}$}
\end{qoptions}
\end{bankq}
\begin{bankq}[id=NUR04-P1-MATH-S1-Q02, source={Nurture (General) - Test 04 (11 Feb 2024)}, paper=1, qno=2, type=single, subtopic={Standard forms}, answer={C}, answer_src=solved, working={Right side is (x-y)/sqrt2, the distance from the line x=y: a parabola with focus (2,1)}, difficulty=6, diff_parts={K2 S2 I2 C1 F0}, diff_note={concepts: locus with an angle condition, standard forms | steps: expand sine; locate focus and directrix; identify the conic}]
If $|z-2-i|=|z|\sin\!\left(\dfrac\pi4-\arg z\right)$, then locus of z is/an:
\begin{qoptions}
\opt{A}{Ellipse}
\opt{B}{Circle}
\opt{C}{Parabola}
\opt{D}{Pair of straight lines}
\end{qoptions}
\end{bankq}
\end{document}
